\section{Sujet n\textsuperscript{o}\,27}
\noindent\begin{minipage}{0.5\linewidth}\footnotesize Office du Baccalauréat du Cameroun\end{minipage}%
\hfill\begin{minipage}{0.45\linewidth}\footnotesize\raggedleft Examen : \textbf{Baccalauréat} · Série : \textbf{C/E}\\
Épreuve : \textbf{Mathématiques} · Durée : \textbf{4\,h} · Coef. : \textbf{7}\end{minipage}
\par\smallskip{\color{onbo}\rule{\linewidth}{1pt}}

\noindent\textit{Un garage automobile de la zone industrielle de Bonabéri (Douala) gère son atelier
de mécanique. Les exercices ci-dessous s'inspirent de cette activité.}

\subsection*{Partie A — Évaluation des ressources \hfill (15 points)}

\noindent\textbf{Exercice 1.} \hfill\textit{(3 points)}\\
Un casier contient $12$ plaquettes de frein dont $8$ sont en bon état et $4$ sont
défectueuses, indiscernables au toucher. Un mécanicien prélève \emph{simultanément} $3$
plaquettes. Soit $X$ le nombre de plaquettes défectueuses prélevées.
\begin{enumerate}[label=\arabic*)]
\item Déterminer la loi de probabilité de $X$. \hfill\textit{(1,75 pt)}
\item Calculer l'espérance mathématique $E(X)$. \hfill\textit{(0,75 pt)}
\item Calculer la probabilité de prélever au moins une plaquette défectueuse. \hfill\textit{(0,5 pt)}
\end{enumerate}

\smallskip
\noindent\textbf{Exercice 2.} \hfill\textit{(3 points)}\\
Le compteur kilométrique d'un véhicule de location affiche un entier. Le garagiste effectue
des calculs modulaires pour son suivi.
\begin{enumerate}[label=\arabic*)]
\item Déterminer le reste de la division euclidienne de $3^{2025}$ par $7$. \hfill\textit{(1,5 pt)}
\item Résoudre dans $\Z$ l'équation $3x\equiv2\ [7]$. \hfill\textit{(1,5 pt)}
\end{enumerate}

\smallskip
\noindent\textbf{Exercice 3.} \hfill\textit{(4 points)}\\
L'usure (en mm) d'un pneu après $x$ milliers de kilomètres est modélisée sur $[0,+\infty[$ par
$f(x)=\dfrac{x}{x+1}+\ln(x+1)$. On note $(\mathcal{C})$ sa courbe.
\begin{enumerate}[label=\arabic*)]
\item Calculer $f(0)$ et $\displaystyle\lim_{x\to+\infty}f(x)$. \hfill\textit{(0,75 pt)}
\item Montrer que $f'(x)=\dfrac{x+2}{(x+1)^2}$, puis dresser le tableau de variations de $f$ sur $[0,+\infty[$. \hfill\textit{(1,5 pt)}
\item Montrer que l'équation $f(x)=2$ admet une unique solution $\alpha$ sur $[0,+\infty[$
(théorème des valeurs intermédiaires), puis justifier l'encadrement $2<\alpha<3$. \hfill\textit{(1,75 pt)}
\end{enumerate}

\smallskip
\noindent\textbf{Exercice 4.} \hfill\textit{(5 points)}\\
L'espace est muni d'un repère orthonormé $(O\,;\,\vec\imath,\vec\jmath,\vec k)$. On considère les
points $A(2\,;\,0\,;\,1)$, $B(1\,;\,2\,;\,0)$, $C(0\,;\,1\,;\,3)$ et $D(3\,;\,2\,;\,4)$.
\begin{enumerate}[label=\arabic*)]
\item Calculer les coordonnées du produit vectoriel $\vec{AB}\wedge\vec{AC}$. \hfill\textit{(1,5 pt)}
\item En déduire une équation cartésienne du plan $(ABC)$. \hfill\textit{(1,25 pt)}
\item Calculer la distance du point $D$ au plan $(ABC)$. \hfill\textit{(1 pt)}
\item En déduire l'aire du triangle $ABC$, puis le volume du tétraèdre $ABCD$. \hfill\textit{(1,25 pt)}
\end{enumerate}

\subsection*{Partie B — Évaluation des compétences \hfill (5 points)}

\noindent\textit{Le garage prépare son année. Le gérant te sollicite, en tant qu'élève de
Terminale~C, pour trois décisions.\\
$\bullet$ Pour stocker l'huile de vidange usagée, il fait fabriquer une cuve cylindrique \emph{ouverte}
(sans couvercle) de volume $V=\SI{1}{\cubic\metre}$. Le coût est proportionnel à la surface de
tôle utilisée (fond $+$ paroi latérale). Il veut minimiser cette surface.\\
$\bullet$ Le nombre de véhicules entretenus était de $800$ en $2024$ et augmente de $50$ véhicules chaque année.\\
$\bullet$ Sur l'ensemble des réparations, \SI{15}{\percent} nécessitent une seconde visite ; on examine $6$ réparations au hasard.}

\begin{enumerate}[label=\textbf{Tâche \arabic*.},leftmargin=2.4em]
\item Déterminer le rayon $r$ et la hauteur $h$ de la cuve cylindrique ouverte de volume
\SI{1}{\cubic\metre} qui minimisent la surface de tôle $S=\pi r^2+2\pi r h$. \hfill\textit{(1,5 pt)}
\item Déterminer l'année à partir de laquelle le nombre de véhicules entretenus dépassera $1200$. \hfill\textit{(1,5 pt)}
\item Déterminer la probabilité qu'\emph{exactement deux} des $6$ réparations examinées nécessitent une seconde visite. \hfill\textit{(1,5 pt)}
\end{enumerate}
\par\smallskip{\footnotesize\itshape Présentation générale : 0,5 point.}

% ════════════════════════════════════════════════════
\corrige{du sujet 27}

\noindent\textbf{Partie A — Exercice 1.}
1) Tirages simultanés de $3$ parmi $12$ : $\binom{12}{3}=220$. Loi hypergéométrique :
$P(X{=}k)=\dfrac{\binom4k\binom{8}{3-k}}{220}$, $k\in\{0,1,2,3\}$.
$P(0)=\frac{\binom40\binom83}{220}=\frac{56}{220}=\frac{14}{55}$ ;
$P(1)=\frac{\binom41\binom82}{220}=\frac{4\cdot28}{220}=\frac{112}{220}=\frac{28}{55}$ ;
$P(2)=\frac{\binom42\binom81}{220}=\frac{6\cdot8}{220}=\frac{48}{220}=\frac{12}{55}$ ;
$P(3)=\frac{\binom43\binom80}{220}=\frac{4}{220}=\frac{1}{55}$.
(Total : $\frac{14+28+12+1}{55}=\frac{55}{55}=1$ $\checkmark$.)
2) $E(X)=0\cdot\frac{14}{55}+1\cdot\frac{28}{55}+2\cdot\frac{12}{55}+3\cdot\frac{1}{55}=\frac{28+24+3}{55}=\frac{55}{55}=1$.
(On retrouve $E(X)=n\frac{N_1}{N}=3\times\frac{4}{12}=1$.)
3) $P(X\ge1)=1-P(0)=1-\frac{14}{55}=\frac{41}{55}\approx0{,}745$.

\noindent\textbf{Exercice 2.}
1) Modulo $7$ : $3^1\equiv3,\ 3^2\equiv2,\ 3^3\equiv6,\ 3^4\equiv4,\ 3^5\equiv5,\ 3^6\equiv1\ [7]$
(ordre $6$). Or $2025=6\times337+3$, donc $3^{2025}\equiv3^{3}\equiv6\ [7]$. Le reste est $\mathbf{6}$.
2) On cherche l'inverse de $3$ modulo $7$ : $3\times5=15\equiv1\ [7]$, donc $3^{-1}\equiv5$.
$3x\equiv2\Rightarrow x\equiv5\times2=10\equiv3\ [7]$. \textbf{Solutions :} $x\equiv3\ [7]$, soit $x=3+7k,\ k\in\Z$.

\noindent\textbf{Exercice 3.}
1) $f(0)=\frac{0}{1}+\ln1=0$. Quand $x\to+\infty$ : $\frac{x}{x+1}\to1$ et $\ln(x+1)\to+\infty$,
donc $\displaystyle\lim_{x\to+\infty}f(x)=+\infty$.
2) $f'(x)=\dfrac{(x+1)-x}{(x+1)^2}+\dfrac{1}{x+1}=\dfrac{1}{(x+1)^2}+\dfrac{1}{x+1}
=\dfrac{1+(x+1)}{(x+1)^2}=\dfrac{x+2}{(x+1)^2}$.
Sur $[0,+\infty[$, $x+2>0$ et $(x+1)^2>0$, donc $f'(x)>0$ : $f$ est strictement croissante,
de $f(0)=0$ vers $+\infty$.
3) $f$ est continue et strictement croissante sur $[0,+\infty[$, de $f(0)=0$ à $+\infty$ ;
comme $2\in[0,+\infty[$, le théorème des valeurs intermédiaires (et la stricte monotonie)
garantit une \emph{unique} solution $\alpha$ à $f(x)=2$.
Encadrement : $f(2)=\frac{2}{3}+\ln3\approx0{,}667+1{,}099=1{,}766<2$ et
$f(3)=\frac{3}{4}+\ln4\approx0{,}75+1{,}386=2{,}136>2$. Comme $f(2)<2<f(3)$ et que $f$ est
strictement croissante, on a $\mathbf{2<\alpha<3}$.

\noindent\textbf{Exercice 4.}
1) $\vec{AB}=B-A=(-1\,;\,2\,;\,-1)$, $\vec{AC}=C-A=(-2\,;\,1\,;\,2)$.
$\vec{AB}\wedge\vec{AC}=\big(2\cdot2-(-1)\cdot1\,;\,(-1)\cdot(-2)-(-1)\cdot2\,;\,(-1)\cdot1-2\cdot(-2)\big)
=(4+1\,;\,2+2\,;\,-1+4)=(5\,;\,4\,;\,3)$.
2) Le plan $(ABC)$ a pour vecteur normal $\vec n=(5\,;\,4\,;\,3)$ et passe par $A(2\,;\,0\,;\,1)$ :
$5(x-2)+4(y-0)+3(z-1)=0\Rightarrow 5x+4y+3z-13=0$.
3) $d(D,(ABC))=\dfrac{|5(3)+4(2)+3(4)-13|}{\sqrt{5^2+4^2+3^2}}=\dfrac{|15+8+12-13|}{\sqrt{50}}
=\dfrac{22}{5\sqrt2}=\dfrac{22\sqrt2}{10}=\dfrac{11\sqrt2}{5}\approx3{,}11$.
4) Aire $ABC=\frac12\|\vec{AB}\wedge\vec{AC}\|=\frac12\sqrt{50}=\frac{5\sqrt2}{2}$.
Volume tétraèdre $=\frac13\times\text{aire}\times d(D,(ABC))=\frac13\times\frac{5\sqrt2}{2}\times\frac{11\sqrt2}{5}
=\frac13\times\frac{5\cdot11\cdot2}{2\cdot5}=\frac13\times11=\dfrac{11}{3}\approx3{,}67$ (unités de volume).

\smallskip
\noindent\textbf{Partie B — Situation.}
\emph{Tâche 1.} Volume $V=\pi r^2 h=1\Rightarrow h=\frac{1}{\pi r^2}$.
Surface (fond $+$ latérale) $S(r)=\pi r^2+2\pi r h=\pi r^2+2\pi r\cdot\frac{1}{\pi r^2}=\pi r^2+\frac{2}{r}$.
$S'(r)=2\pi r-\frac{2}{r^2}=\frac{2\pi r^3-2}{r^2}$. $S'(r)=0\Rightarrow r^3=\frac{1}{\pi}\Rightarrow
r=\big(\tfrac1\pi\big)^{1/3}\approx\SI{0,683}{\metre}$.
$S''(r)=2\pi+\frac{4}{r^3}>0$ : minimum. Alors $h=\frac{1}{\pi r^2}=\frac{1}{\pi}\cdot\pi^{2/3}=\pi^{-1/3}=r$,
donc $h=r\approx\SI{0,68}{\metre}$.

\emph{Tâche 2.} Suite arithmétique de raison $50$, premier terme (en $2024$, $n=0$) : $v_n=800+50n$.
$v_n>1200\Rightarrow 50n>400\Rightarrow n>8$, donc à partir de $n=9$, soit l'année \textbf{2033}
($v_9=800+450=1250$).

\emph{Tâche 3.} Le nombre de réparations nécessitant une seconde visite suit $\mathcal{B}(6\,;\,0{,}15)$.
$P(X{=}2)=\binom62(0{,}15)^2(0{,}85)^4=15\times0{,}0225\times0{,}52200625\approx\mathbf{0{,}176}$,
soit environ \SI{17,6}{\percent}.
